\chapter{Introduction}
\label{ch:introduction}

\section{Context and Motivation}
\label{sec:intro-context}

Whistleblowers occupy a uniquely fragile position in modern society. They are
the individuals who, from inside an organisation, decide that the public
interest outweighs their personal safety, and who choose to expose corruption,
fraud, illegal surveillance, environmental damage, or abuse. The historical
record is unambiguous about their importance: many of the defining
accountability stories of the past two decades---from mass-surveillance
disclosures to corporate financial scandals---began with a single person who
held evidence and decided to share it. It is equally unambiguous about the price
they pay. Whistleblowers routinely lose their jobs, face civil and criminal
litigation, are placed under surveillance, and in the most severe cases are
detained or physically threatened.

The technical means available to protect such a person have not kept pace with
the threats they face. The dominant model for handling sensitive disclosures is
still \emph{custodial}: the source hands their material to a trusted
intermediary---a newsroom, a non-governmental organisation, or a dedicated
submission platform---and relies on that intermediary to safeguard both the data
and their anonymity. This model concentrates risk in exactly the place an
adversary will attack. A custodian can be subpoenaed, infiltrated, hacked, or
simply coerced; once it is compromised, every source who ever trusted it is
exposed at once. Centralised submission systems such as SecureDrop mitigate some
of these risks through careful operational security, but they remain
single points of organisational trust and they do nothing to solve two problems
that sit at the heart of high-stakes whistleblowing.

The first problem is the \emph{vulnerable interval}. Sensitive evidence is at
its most dangerous to its holder precisely while it is unpublished. During that
interval, an adversary has a clear incentive to act: if the source can be
silenced---through intimidation, detention, or worse---before the material is
released, the disclosure never happens. The source's possession of the evidence,
which ought to be their protection, becomes the very thing that endangers them.

The second problem is \emph{credible anonymity}. A source who wishes to remain
anonymous cannot prove who they are, and a journalist receiving an anonymous tip
has no way to distinguish a genuine insider from a fabricator or a deliberate
disinformation campaign. The conventional resolution---revealing identity to
establish credibility---defeats the purpose of anonymity entirely. Compounding
this, when a tip turns out to be valuable, there is no safe way to reward the
source: any payment creates a financial trail that can be followed straight back
to them.

This thesis is motivated by the conviction that these problems are not
fundamentally social or legal but \emph{architectural}, and that they can be
addressed with cryptography and decentralised infrastructure that have only
recently become mature enough to use in production. The work presented here
designs, implements, and evaluates such a system.

\section{The Problem Statement}
\label{sec:intro-problem}

The concept of a \emph{dead man's switch} (DMS) is a direct answer to the
vulnerable-interval problem. A dead man's switch is a mechanism that performs an
action automatically unless it is periodically and deliberately prevented from
doing so. Applied to whistleblowing, the action is the release of encrypted
evidence, and the prevention is a regular ``check-in'' by the source. As long as
the source is free and safe, they check in and the evidence stays sealed. If
they are detained, incapacitated, or killed, the check-ins stop, and the
evidence is released automatically. The switch inverts the adversary's
incentive: silencing the source no longer suppresses the disclosure---it
\emph{triggers} it.

The difficulty has always been \emph{who or what holds the switch}. A human
confidant can be coerced or can lose their nerve. A centralised online service
can be shut down, subpoenaed, or can simply fail. Any custodian that is powerful
enough to release the data on the source's behalf is also powerful enough to
release it prematurely, withhold it indefinitely, or be compelled to hand over
the keys. The release mechanism must therefore be \emph{autonomous} (it runs
without a trusted operator), \emph{unstoppable} (no single party can prevent the
release), and \emph{confidential} (no single party can perform the release
early).

An earlier iteration of the project on which this thesis is based attempted to
solve the custody problem using a \emph{Trusted Execution Environment} (TEE)---a
hardware enclave, specifically Intel SGX accessed through the iExec network, in
which decryption keys could be sealed and used without the host operator seeing
them. As Chapter~\ref{ch:background} discusses in detail, this approach trades a
human single point of failure for a \emph{hardware} single point of failure. It
relies on the integrity of a specific chip vendor and on the assumption that the
enclave cannot be broken into---an assumption repeatedly undermined by
side-channel and transient-execution attacks. For an adversary with
nation-state resources, a hardware root of trust is not a guarantee; it is a
target.

The problem this thesis addresses can therefore be stated precisely:

\begin{quote}
\emph{How can a whistleblower seal sensitive evidence such that it is released
automatically if and only if they stop checking in, without entrusting the
decryption capability, the storage, or their identity and rewards to any single
party---hardware, human, or organisation---that could be coerced, compromised,
or compelled?}
\end{quote}

A complete answer must simultaneously address four sub-problems: (i) where the
encrypted evidence is stored so that it cannot be censored or deleted; (ii) who
holds the decryption key and under what condition they release it; (iii) how an
anonymous source can nonetheless establish credibility; and (iv) how a source
can be paid without being deanonymised.

\section{Proposed Solution}
\label{sec:intro-solution}

This thesis presents \zkwns, a decentralised, censorship-resistant
whistleblowing platform that answers each of the four sub-problems with a
distinct, production-ready cryptographic primitive, and integrates them into one
coherent system. The guiding design principle is the deliberate replacement of
\emph{hardware-rooted} and \emph{custodial} trust with trust rooted in
\emph{mathematics} and \emph{decentralised consensus}---a move the project
characterises throughout as ``from hardware to math''.

The system is organised into three functional modules:

\begin{itemize}
  \item \textbf{Module A --- The Vault (Dead Man's Switch).} A file is encrypted
  in the source's browser with AES-256-GCM; the ciphertext is stored permanently
  on Arweave; and the symmetric key is sealed by the \emph{Lit Protocol}, a
  decentralised key-management network that uses multi-party computation and
  threshold cryptography. The key is released only when an on-chain
  smart-contract condition, \isdeceased, evaluates to true---which happens
  exactly when the source misses their check-in interval.

  \item \textbf{Module B --- Identity and Reputation (Zero-Knowledge
  Provenance).} Using the \emph{Reclaim Protocol} and zkTLS, the source generates
  a zero-knowledge proof of a private HTTPS session---for example, that they are
  an active employee of a named organisation---and anchors a hash of that proof
  on-chain. A journalist can thereby verify the source's credibility without ever
  learning who they are.

  \item \textbf{Module C --- The Marketplace (Anonymous Exchange).} Sources list
  encrypted information and journalists place escrowed bids. When a bid is
  accepted, payment is routed to an \emph{ERC-5564 stealth address}, severing the
  on-chain link between the payer and the payee.
\end{itemize}

Crucially, \emph{no plaintext and no key material ever leaves the source's
device in the clear}, and \emph{no single component of the system can both access
the ciphertext and release the key} ahead of the agreed condition. Each layer is
operated by a different decentralised network with a different failure mode, so
that compromising the whole requires compromising several independent systems
simultaneously.

\section{Objectives}
\label{sec:intro-objectives}

The principal objectives of this thesis are the following:

\begin{enumerate}
  \item \textbf{O1 --- Eliminate the hardware trust anchor.} Replace the
  TEE-based key custody of the earlier design with a key-management scheme whose
  security rests on threshold cryptography, and justify the change with a
  rigorous comparison of the two trust models.

  \item \textbf{O2 --- Build an autonomous, unstoppable release mechanism.}
  Design and implement an on-chain dead man's switch whose trigger condition can
  be read trustlessly by the key-management network, so that release requires no
  human operator.

  \item \textbf{O3 --- Provide credible anonymity.} Integrate zero-knowledge
  provenance so that a source can prove the credibility of a claim without
  revealing their identity, and anchor that proof on-chain in a replay-resistant
  way.

  \item \textbf{O4 --- Enable unlinkable rewards.} Implement anonymous payments
  via stealth addresses so that valuable sources can be compensated without
  creating a traceable financial link.

  \item \textbf{O5 --- Guarantee data permanence and client-side
  confidentiality.} Ensure all encryption happens on the client and that the
  encrypted payload is stored on censorship-resistant permanent storage.

  \item \textbf{O6 --- Validate the system.} Implement the design as audited-style
  smart contracts and a typed client, accompany it with an automated test suite,
  and deploy it to a public test network to demonstrate an end-to-end flow.
\end{enumerate}

\section{Contributions}
\label{sec:intro-contributions}

The main contributions of this work are:

\begin{enumerate}
  \item A \textbf{critical re-evaluation of the trust model} for confidential
  whistleblowing, articulating precisely why a hardware-enforced enclave is an
  inappropriate root of trust in a nation-state adversary model and why threshold
  cryptography is a better fit (Chapters~\ref{ch:background}
  and~\ref{ch:evaluation}). This analysis is sharpened by a discussion of the Lit
  Protocol's own partial return to TEE-based execution, and a justification of
  the design's deliberate choice to remain on the multi-party-computation path.

  \item A \textbf{coherent four-primitive architecture} that combines a
  cryptographic dead man's switch, threshold key management, zero-knowledge
  provenance, and stealth-address payments into a single end-to-end system, with
  a clear specification of what is and is not stored on-chain
  (Chapter~\ref{ch:architecture}).

  \item A \textbf{complete reference implementation} comprising three Solidity
  smart contracts (a heartbeat-based switch, a reputation registry with on-chain
  zkTLS proof verification, and an escrowed marketplace with stealth payouts),
  five client-side cryptographic services, hardened server routes that keep
  application secrets off the client, and a multi-step user interface
  (Chapter~\ref{ch:implementation}).

  \item A \textbf{security-engineering treatment} of the smart contracts and
  services---layered verification trust modes, the checks-effects-interactions
  pattern reinforced with reentrancy guards, replay-resistant proof binding, and
  server-side secret handling---together with an automated test suite of
  \textbf{91 passing tests} and a public deployment to the Base Sepolia test
  network (Chapters~\ref{ch:implementation} and~\ref{ch:evaluation}).

  \item An \textbf{honest evaluation} against the stated objectives and threat
  model, including a frank account of the prototype's current limitations and a
  concrete roadmap of future work (Chapter~\ref{ch:evaluation} and
  Chapter~\ref{ch:conclusions}).
\end{enumerate}

\section{Thesis Outline}
\label{sec:intro-outline}

The remainder of this thesis is organised as follows.

\textbf{Chapter~\ref{ch:background}, Background and Theoretical Foundations},
introduces the cryptographic and blockchain concepts on which the system is
built---authenticated symmetric encryption, elliptic-curve cryptography,
threshold cryptography and multi-party computation, the EVM and smart contracts,
zero-knowledge proofs and zkTLS, stealth addresses, and permanent decentralised
storage---and explains the limitations of trusted execution environments that
motivate the central design decision.

\textbf{Chapter~\ref{ch:related}, State of the Art}, surveys existing
whistleblowing platforms, dead-man's-switch implementations, decentralised
key-management networks, and provenance technologies, and positions \zkw within
this landscape through a gap analysis.

\textbf{Chapter~\ref{ch:architecture}, System Design and Architecture}, presents
the design goals and threat model, the layered architecture, the rationale for
the move from hardware to mathematics, the design of each of the three modules,
the data model, and the end-to-end workflows.

\textbf{Chapter~\ref{ch:implementation}, Implementation}, documents the concrete
realisation of the design: the technology stack, the three smart contracts and
their deployment, the client-side cryptographic services, the integration with
Lit, Reclaim, Arweave and stealth addresses, the hardened server routes, the
user interface, and the security-engineering decisions throughout.

\textbf{Chapter~\ref{ch:evaluation}, Testing, Deployment and Evaluation},
describes the testing methodology and results, the public testnet deployment, the
end-to-end validation, an evaluation against the objectives and threat model, and
a candid discussion of limitations.

\textbf{Chapter~\ref{ch:conclusions}, Conclusions and Future Work}, summarises
the contributions, reflects on the broader significance of the design, and
outlines avenues for future development.
